\chapter{第6章：Jacobian 习题}
\difficultykey

\section{基础题 \easy}

\begin{problem}{模曲线 Jacobian 的维数}
\easy
利用
\[
\dim J_0(N)=g\bigl(X_0(N)\bigr)=\dim S_2(\Gamma_0(N))
\]
计算 $N=11,17,19,23,37$ 时 $J_0(N)$ 的维数。
\end{problem}
\begin{solution}
由第 5 章的权 $2$ 维数表可知
\[
\dim S_2(\Gamma_0(11))=
\dim S_2(\Gamma_0(17))=
\dim S_2(\Gamma_0(19))=1,
\]
以及
\[
\dim S_2(\Gamma_0(23))=
\dim S_2(\Gamma_0(37))=2.
\]
所以
\[
\dim J_0(11)=\dim J_0(17)=\dim J_0(19)=1,
\]
\[
\dim J_0(23)=\dim J_0(37)=2.
\]
前面三个 Jacobian 是椭圆曲线，后面两个是二维 Abel 簇。
\end{solution}

\begin{problem}{基点改变只差平移}
\easy
设 $X$ 是紧黎曼面，$P_0,Q_0\in X$。证明以 $P_0$ 与 $Q_0$ 为基点定义的 Abel--Jacobi 映射只相差 $J(X)$ 上的一个平移。
\end{problem}
\begin{solution}
记两条 Abel--Jacobi 映射分别为 $\phi_{P_0}$ 与 $\phi_{Q_0}$。对任意 $P\in X$ 和任意全纯微分 $\omega$，有
\[
\phi_{Q_0}(P)(\omega)-\phi_{P_0}(P)(\omega)=\int_{Q_0}^{P}\omega-\int_{P_0}^{P}\omega=\int_{Q_0}^{P_0}\omega.
\]
右边与点 $P$ 无关，只取决于两个基点。所以存在一个固定元素 $c\in J(X)$ 使
\[
\phi_{Q_0}(P)=\phi_{P_0}(P)+c
\]
对所有 $P$ 都成立。
这正说明两者只差一个平移。
\end{solution}

\begin{problem}{次数零除子类}
\easy
解释为什么 $J(X)$ 也可以看成 $\operatorname{Pic}^0(X)$，也就是次数为零的除子类群。
\end{problem}
\begin{solution}
解析定义中，$J(X)$ 是全纯微分对偶空间模掉周期格：
\[
J(X)=H^0(X,\Omega^1)^*/H_1(X,\Z).
\]
Abel 定理与 Jacobi 逆问题告诉我们：一个次数零除子
\[
D=\sum_i n_i(P_i)
\qquad \left(\sum_i n_i=0\right)
\]
的类只由积分
\[
\omega\longmapsto \sum_i n_i\int_{P_0}^{P_i}\omega
\]
决定，而这个积分恰好只在 Jacobian 中定义到周期为止。

因此“次数零除子模主除子”与“积分值模周期格”给出同一个对象，于是
\[
J(X)\cong \operatorname{Pic}^0(X).
\]
\end{solution}

\begin{problem}{亏格一曲线等于自己的 Jacobian}
\easy
设 $X$ 是亏格 $1$ 的紧黎曼面，并固定一点 $P_0\in X$。说明 Abel--Jacobi 映射
\[
\phi_{P_0}:X\to J(X)
\]
是同构。
\end{problem}
\begin{solution}
这是亏格 $1$ 情形的基本定理。因为 $g(X)=1$，全纯微分空间是一维，周期格在其对偶空间中给出一个秩 $2$ 的格，所以
\[
J(X)\cong \C/\Lambda
\]
是一条椭圆曲线。

另一方面，亏格 $1$ 的紧黎曼面本身也可写成某个复环面 $\C/\Lambda'$。Abel--Jacobi 映射由积分构造，正好把 $X$ 与它自己的周期环面识别起来，所以是双全纯同构。

因此亏格 $1$ 曲线的 Jacobian 就是它自己。
\end{solution}

\begin{problem}{一级别十一的 Jacobian 是椭圆曲线}
\easy
说明 $J_0(11)$ 是一条椭圆曲线。
\end{problem}
\begin{solution}
由第 5 章，
\[
g\bigl(X_0(11)\bigr)=1.
\]
因此
\[
\dim J_0(11)=g\bigl(X_0(11)\bigr)=1.
\]
一维 Abel 簇就是椭圆曲线，所以 $J_0(11)$ 是一条椭圆曲线。

也可以这样理解：因为 $X_0(11)$ 本身是亏格 $1$ 曲线，所以它与自己的 Jacobian 同构，而 Jacobian 必是一维 Abel 簇。
\end{solution}

\begin{problem}{二挠点的大小}
\easy
设 $J$ 是维数 $g$ 的复 Abel 簇。证明 $J[2]\cong (\Z/2\Z)^{2g}$，并据此求出 $J_0(11)[2]$ 的大小。
\end{problem}
\begin{solution}
把 $J$ 写成复环面 $V/\Lambda$，其中 $\Lambda\cong \Z^{2g}$。一个点 $v+\Lambda$ 满足 $2(v+\Lambda)=0$ 当且仅当 $2v\in\Lambda$，也就是
\[
v\in \frac{1}{2}\Lambda.
\]
所以
\[
J[2]\cong \frac{\frac{1}{2}\Lambda}{\Lambda}\cong \Lambda/2\Lambda\cong (\Z/2\Z)^{2g}.
\]

对 $J_0(11)$，我们知道 $g=1$，因此
\[
J_0(11)[2]\cong (\Z/2\Z)^2.
\]
所以它恰有
\[
2^2=4
\]
个二挠点。
\end{solution}

\begin{problem}{一般挠子群与 Tate 模的秩}
\easy
设 $J$ 是维数 $g$ 的 Abel 簇，$\ell$ 为素数。说明
\[
J[\ell^n]\cong (\Z/\ell^n\Z)^{2g}
\]
并推出 Tate 模 $T_\ell(J)$ 是自由的 $\Z_\ell$-模，秩为 $2g$。
\end{problem}
\begin{solution}
仍把 $J$ 写成 $V/\Lambda$，其中 $\Lambda\cong \Z^{2g}$。与上一题完全同样的计算给出
\[
J[\ell^n]\cong \Lambda/\ell^n\Lambda\cong (\Z/\ell^n\Z)^{2g}.
\]

定义
\[
T_\ell(J)=\varprojlim_n J[\ell^n].
\]
对逆极限取上面的同构，得到
\[
T_\ell(J)\cong \varprojlim_n (\Z/\ell^n\Z)^{2g}\cong \Z_\ell^{2g}.
\]
因此 $T_\ell(J)$ 是自由的 $\Z_\ell$-模，秩恰为 $2g$。
\end{solution}

\begin{problem}{两个具体 Tate 模的秩}
\easy
计算 $T_\ell(J_0(11))$ 与 $T_\ell(J_0(37))$ 的 $\Z_\ell$-秩。
\end{problem}
\begin{solution}
一般结论是：若 $\dim J=g$，则 $T_\ell(J)$ 的秩为 $2g$。

对 $J_0(11)$，有 $g=1$，所以
\[
T_\ell(J_0(11))\cong \Z_\ell^2.
\]
对 $J_0(37)$，有 $g=2$，所以
\[
T_\ell(J_0(37))\cong \Z_\ell^4.
\]
因此它们的秩分别是 $2$ 与 $4$。
\end{solution}

\begin{problem}{一级别十一上的 Hecke 作用}
\easy
设 $f_{11}=q-2q^2-q^3+2q^4+q^5+\cdots$ 是 $S_2(\Gamma_0(11))$ 的唯一归一化新形式。求 $T_2$ 与 $T_3$ 在 $H^0(J_0(11),\Omega^1)$ 上的作用标量。
\end{problem}
\begin{solution}
由 $H^0(J_0(11),\Omega^1)\cong S_2(\Gamma_0(11))$，Hecke 算子在余切空间上的作用就是经典 Hecke 作用。因此
\[
T_n\text{ 在 }H^0(J_0(11),\Omega^1)\text{ 上按 }a_n(f_{11})\text{ 作用}.
\]
从 $q$-展开读出
\[
a_2(f_{11})=-2,
\qquad
 a_3(f_{11})=-1.
\]
所以
\[
T_2=-2,
\qquad
T_3=-1
\]
作为这一维空间上的作用标量。
\end{solution}

\section{进阶题 \medium}

\begin{problem}{周期格为什么真是格}
\medium
设 $X$ 是亏格 $g$ 的紧黎曼面。解释为什么周期映射
\[
H_1(X,\Z)\longrightarrow H^0(X,\Omega^1)^*
\]
的像是一个秩 $2g$ 的离散子群。
\end{problem}
\begin{solution}
首先，$H_1(X,\Z)$ 是秩 $2g$ 的自由阿贝尔群。其次，$H^0(X,\Omega^1)$ 的复维数是 $g$，因此其对偶空间的实维数是 $2g$。

周期映射把一个同调类 $\gamma$ 送到线性泛函
\[
\omega\longmapsto \int_\gamma \omega.
\]
由 de Rham 配对的非退化性，这个映射在实向量空间层面是单射。由于定义域与值域的实维数同为 $2g$，它事实上给出一个实线性同构
\[
H_1(X,\R)\cong H^0(X,\Omega^1)^*.
\]
于是整数格 $H_1(X,\Z)$ 的像必然是对偶空间中的一个离散满秩子群，也就是秩 $2g$ 的格。
\end{solution}

\begin{problem}{Abel--Jacobi 对除子的可加性}
\medium
固定基点 $P_0\in X$。设
\[
D=\sum_i n_i(P_i)
\qquad \left(\sum_i n_i=0\right)
\]
是一个次数零除子。说明元素
\[
\sum_i n_i\,\phi_{P_0}(P_i)\in J(X)
\]
只依赖于 $D$，并且对除子加法是可加的。
\end{problem}
\begin{solution}
由 Abel--Jacobi 映射定义，$\phi_{P_0}(P_i)$ 是从 $P_0$ 积分到 $P_i$ 得到的周期类。于是
\[
\sum_i n_i\phi_{P_0}(P_i)
\]
对应的线性泛函是
\[
\omega\longmapsto \sum_i n_i\int_{P_0}^{P_i}\omega.
\]
因为 $\sum_i n_i=0$，若把每条路径都改动一个公共闭回路，只会增加总系数为零的周期和，所以该类良定义。

它显然只依赖于除子系数与点的位置，因此只依赖于 $D$。又因为积分与求和都是线性的，所以
\[
\Phi(D_1+D_2)=\Phi(D_1)+\Phi(D_2).
\]
这就是可加性。
\end{solution}

\begin{problem}{同态的线性提升与有限核}
\medium
设 $A_i=V_i/\Lambda_i$ 是两个同维数的复环面。说明：若复线性映射 $L:V_1\to V_2$ 满足 $L(\Lambda_1)\subset \Lambda_2$，则它诱导同态 $\varphi:A_1\to A_2$；若 $L$ 还是线性同构且 $L(\Lambda_1)$ 在 $\Lambda_2$ 中有限指数，则 $\varphi$ 有限核且为满射。
\end{problem}
\begin{solution}
首先定义
\[
\varphi(v+\Lambda_1)=L(v)+\Lambda_2.
\]
若 $v$ 改成 $v+\lambda$，其中 $\lambda\in\Lambda_1$，则
\[
L(v+\lambda)-L(v)=L(\lambda)\in\Lambda_2,
\]
所以定义与代表元无关，故得到复环面同态。

若 $L$ 是线性同构，则每个 $w\in V_2$ 都写成 $w=L(v)$，所以 $\varphi$ 满射。再看核：
\[
\ker\varphi=\{v+\Lambda_1: L(v)\in \Lambda_2\}.
\]
由 $L$ 可逆，核与商群 $\Lambda_2/L(\Lambda_1)$ 同构。若后者有限，则核有限。
因此 $\varphi$ 是同种。
\end{solution}

\begin{problem}{乘法映射的次数}
\medium
设 $A=V/\Lambda$ 是维数 $g$ 的 Abel 簇。证明乘法映射
\[
[m]:A\to A,
\qquad x\longmapsto mx
\]
的核同构于 $(\Z/m\Z)^{2g}$，因此其次数为 $m^{2g}$。
\end{problem}
\begin{solution}
点 $v+\Lambda$ 落在核中，当且仅当
\[
mv\in \Lambda.
\]
也就是
\[
v\in \frac{1}{m}\Lambda.
\]
因此
\[
\ker[m]\cong \frac{\frac{1}{m}\Lambda}{\Lambda}\cong \Lambda/m\Lambda\cong (\Z/m\Z)^{2g}.
\]
核共有 $m^{2g}$ 个元素。

对 Abel 簇同态，有限核的大小就是同种的次数，所以
\[
\deg[m]=m^{2g}.
\]
这推广了一维椭圆曲线情形的熟悉公式。
\end{solution}

\begin{problem}{Hecke 对应给出 Jacobian 自同态}
\medium
设
\[
X_0(N) \xleftarrow{\ \alpha_n\ } X_0(N;n) \xrightarrow{\ \beta_n\ } X_0(N)
\]
是 Hecke 对应。说明组合
\[
T_n=(\beta_n)_*\circ \alpha_n^*
\]
在 $\operatorname{Pic}^0(X_0(N))$ 上良定义，因此给出 $J_0(N)$ 的一个自同态。
\end{problem}
\begin{solution}
先看次数。若 $D$ 是次数零除子，则拉回 $\alpha_n^*D$ 的次数仍为零；再把它推出到 $X_0(N)$ 上，次数仍保持为零。因此 $T_n$ 把次数零除子送到次数零除子。

再看主除子。若 $D=(f)$ 是主除子，则
\[
\alpha_n^*(D)=(f\circ \alpha_n)
\]
仍是主除子。有限映射下推出主除子等于范数函数的除子，所以
\[
(\beta_n)_*\alpha_n^*(D)
\]
仍是主除子。

因此 $T_n$ 在“次数零除子模主除子”的商上良定义，也就是在
\[
\operatorname{Pic}^0(X_0(N))\cong J_0(N)
\]
上良定义。
\end{solution}

\begin{problem}{Hecke 作用与微分一致}
\medium
说明 Jacobian 上由对应得到的 $T_n$，在余切空间 $H^0(J_0(N),\Omega^1)$ 上的作用，就是 $S_2(\Gamma_0(N))$ 上的经典 Hecke 算子。
\end{problem}
\begin{solution}
Jacobian 上的 Hecke 算子是“先拉回除子，再推出除子”。把这个过程在线性化后，余切空间上的作用变成“先拉回微分，再取迹推出微分”。也就是说，在微分空间中得到的是
\[
H^0(X_0(N),\Omega^1)\xrightarrow{\alpha_n^*} H^0(X_0(N;n),\Omega^1)
\xrightarrow{(\beta_n)_*} H^0(X_0(N),\Omega^1).
\]

而经典 Hecke 算子的模解释恰好就是这个“拉回再推出”的过程。再利用同构
\[
H^0(X_0(N),\Omega^1)\cong S_2(\Gamma_0(N)),
\]
便得到两种 Hecke 作用完全一致。
\end{solution}

\begin{problem}{Weil 配对的基本后果}
\medium
设 $J$ 是维数 $g$ 的 Abel 簇。已知 Weil 配对给出一个完美交替配对
\[
T_\ell(J)\times T_\ell(J)\to \Z_\ell(1).
\]
说明这会迫使 $T_\ell(J)$ 的秩是偶数，并解释为什么这与秩 $2g$ 相符。
\end{problem}
\begin{solution}
在线性代数里，一个自由 $\Z_\ell$-模若带有完美交替配对，则它的秩必为偶数；因为交替矩阵在任何基下都可化成若干个
\[
\begin{pmatrix}
0&1\\
-1&0
\end{pmatrix}
\]
块的直和。

所以如果 $T_\ell(J)$ 带有完美交替配对，它的秩只能是偶数。另一方面，我们已经知道
\[
\operatorname{rank}_{\Z_\ell} T_\ell(J)=2g.
\]
这恰好是偶数，与 Weil 配对完全一致。几何上，这说明 Tate 模天然是一个 $2g$ 维辛模。
\end{solution}

\begin{problem}{模 Abel 簇的维数}
\medium
设 $f$ 是权 $2$ 归一化新形式，系数域为
\[
K_f=\Q(a_n(f):n\ge 1).
\]
说明对应模 Abel 簇 $A_f$ 的维数等于 $[K_f:\Q]$，并据此分别求出当 $K_f=\Q$、$[K_f:\Q]=2$、$[K_f:\Q]=3$ 时 $\dim A_f$。
\end{problem}
\begin{solution}
按照定义，$A_f$ 是从 $J_0(N)$ 中切下来的 $f$-方向商。它的余切空间由 $f$ 的全部 Galois 共轭张成，因此维数正好等于这些共轭的个数，也就是系数域次数：
\[
\dim A_f=[K_f:\Q].
\]
于是：
\[
K_f=\Q \Rightarrow \dim A_f=1,
\]
\[
[K_f:\Q]=2 \Rightarrow \dim A_f=2,
\]
\[
[K_f:\Q]=3 \Rightarrow \dim A_f=3.
\]
特别地，系数域有理时，$A_f$ 就是一条椭圆曲线。
\end{solution}

\begin{problem}{一维权二空间时 Jacobian 就是模 Abel 簇}
\medium
设 $\dim S_2(\Gamma_0(N))=1$，记唯一归一化新形式为 $f$。证明 $A_f\cong J_0(N)$。并把结论应用到 $N=11,17,19$。
\end{problem}
\begin{solution}
因为 $S_2(\Gamma_0(N))$ 只有一维，整个余切空间
\[
H^0(J_0(N),\Omega^1)\cong S_2(\Gamma_0(N))
\]
都由 $f$ 这一条 Hecke 方向生成。于是从 $J_0(N)$ 切出 $f$-方向的商时，不会丢掉任何微分方向，所以所得商仍是一维且与原 Jacobian 同构：
\[
A_f\cong J_0(N).
\]

对 $N=11,17,19$，第 5 章已经算出
\[
\dim S_2(\Gamma_0(N))=1.
\]
故这三个级别都满足
\[
A_f\cong J_0(N).
\]
特别地，它们都是椭圆曲线。
\end{solution}

\begin{problem}{二维 Jacobian 的分解图景}
\medium
设 $S_2(\Gamma_0(37))$ 由两条有理新形式生成。说明这意味着 $J_0(37)$ 与两条椭圆曲线的乘积同种。
\end{problem}
\begin{solution}
因为 $S_2(\Gamma_0(37))$ 的两条新形式都有有理 Fourier 系数，所以它们各自的系数域都是 $\Q$。由上一题，附着的模 Abel 簇都一维，也就是椭圆曲线，记为 $E_1,E_2$。

Hecke 理论把 Jacobian 分解成各个新形式方向对应的 Abel 簇的乘积（在同种意义下），所以
\[
J_0(37)\sim E_1\times E_2.
\]
这里“$\sim$”表示同种。由于两边维数都是 $2$，这正好解释了为什么 $J_0(37)$ 是二维 Jacobian，却又可以分解成两个一维模因子。
\end{solution}

\section{挑战题 \hard}

\begin{problem}{亏格一时 Abel--Jacobi 为何是同构}
\hard
给出一个较完整的论证：若 $X$ 是亏格 $1$ 的紧黎曼面，则 Abel--Jacobi 映射
\[
\phi_{P_0}:X\to J(X)
\]
是双全纯同构。
\end{problem}
\begin{solution}
取 $0\ne \omega\in H^0(X,\Omega^1)$。因为 $g=1$，规范除子的次数是
\[
2g-2=0.
\]
全纯微分零点的总次数等于这个数，所以 $\omega$ 没有零点。

设 $\widetilde X\to X$ 是普遍覆盖。定义
\[
F(Q)=\int_{\widetilde P_0}^{Q}\widetilde\omega.
\]
由于 $\widetilde X$ 单连通，这个积分良定义，且
\[
dF=\widetilde\omega.
\]
因为 $\widetilde\omega$ 处处非零，$F$ 是局部双全纯映射。其像是 $\C$ 中既开又闭的连通子集，只能是整个 $\C$，所以 $\widetilde X\cong \C$。

覆盖变换对 $F$ 的作用只能是平移，而这些平移向量正是周期格 $\Lambda_X$。于是
\[
X\cong \C/\Lambda_X.
\]
另一方面，因 $H^0(X,\Omega^1)$ 一维，其对偶空间可与 $\C$ 识别，Jacobian 正是
\[
J(X)=\C/\Lambda_X.
\]
在这个识别下，Abel--Jacobi 映射正是商映射诱导的同构，因此它是双全纯同构。
\end{solution}

\begin{problem}{Galois 共轭张成余切空间}
\hard
设 $f$ 是权 $2$ 归一化新形式，系数域为 $K_f$。说明 $H^0(A_f,\Omega^1)$ 由所有共轭形式 $f^\sigma$ 张成，并重新推出
\[
\dim A_f=[K_f:\Q].
\]
\end{problem}
\begin{solution}
定义 $A_f$ 时，我们把 Hecke 理想 $I_f$ 对应的其余方向全部取商掉，因此 $H^0(A_f,\Omega^1)$ 恰好是 $S_2(\Gamma_0(N))$ 中满足
\[
Th=\lambda_f(T)h
\qquad (T\in \mathbb T)
\]
这一组特征值条件的子空间。

若 $\sigma:K_f\hookrightarrow \C$ 是一个嵌入，则
\[
T(f^\sigma)=\sigma(a_T(f))f^\sigma.
\]
特别地，对属于 $I_f$ 的 Hecke 元，右边为零，所以每个 $f^\sigma$ 都落在这个余切空间里。

反过来，若某个特征方向也被 $I_f$ 消掉，那么它的全部 Hecke 特征值必须与 $f$ 的特征值系统一致，于是只能来自某个 Galois 共轭。故
\[
H^0(A_f,\Omega^1)=\operatorname{span}_\C\{f^\sigma\}.
\]
不同嵌入给出不同的 $q$-展开，因此这些共轭线性无关。它们的个数正是 $[K_f:\Q]$，所以
\[
\dim H^0(A_f,\Omega^1)=[K_f:\Q].
\]
而 Abel 簇维数等于其全纯微分空间维数，于是
\[
\dim A_f=[K_f:\Q].
\]
\end{solution}

\begin{problem}{一级别十一与显式椭圆曲线模型}
\hard
设已知 $X_0(11)$ 可由椭圆曲线模型
\[
E: y^2+y=x^3-x^2-10x-20
\]
给出，并把无穷远点取作基点。说明 $J_0(11)\cong E$。
\end{problem}
\begin{solution}
由第 5 章可知 $g(X_0(11))=1$，所以 $X_0(11)$ 是亏格 $1$ 曲线。既然题目已经给出它的一个椭圆曲线模型 $E$，这就是说
\[
X_0(11)\cong E
\]
作为带基点的亏格 $1$ 曲线。

另一方面，亏格 $1$ 曲线与其 Jacobian 同构：
\[
X_0(11)\cong J(X_0(11))=J_0(11).
\]
把这两个同构复合起来，就得到
\[
J_0(11)\cong E.
\]
也就是说，$J_0(11)$ 不只是抽象的一维 Abel 簇，而且正可以由这条显式椭圆曲线来表示。
\end{solution}

\begin{problem}{模 Abel 簇的泛性质}
\hard
设 $f$ 是权 $2$ 新形式，$q_f:J_0(N)\twoheadrightarrow A_f$ 是对应的商映射。证明：若 $u:J_0(N)\to B$ 是一个 Abel 簇同态，并且对所有 Hecke 算子都满足
\[
u\circ T=\lambda_f(T)\,u,
\]
则 $u$ 必经由 $A_f$ 唯一分解。
\end{problem}
\begin{solution}
由条件可知，对任意 $T\in I_f=\ker\lambda_f$，都有
\[
u\circ T=0.
\]
因此由所有 $T\in I_f$ 的像生成的 Abel 子簇 $I_fJ_0(N)$ 都落在 $\ker u$ 中。

但 $A_f$ 正是把这个子簇压掉得到的商：
\[
A_f=J_0(N)/I_fJ_0(N).
\]
于是商映射的泛性质告诉我们，存在唯一同态
\[
\widetilde u:A_f\to B
\]
使得
\[
u=\widetilde u\circ q_f.
\]
这就是说，$A_f$ 是“承载特征值系统 $\lambda_f$ 的最大商”。
\end{solution}

\begin{problem}{固定向量与有理挠点}
\hard
设 $J/\Q$ 是 Abel 簇，$\ell$ 为素数。说明如果 $T_\ell(J)$ 中存在非零的 Galois 不变向量，那么对每个 $n\ge 1$，$J(\Q)$ 中都会存在一个有理的 $\ell^n$-挠点。由此解释：若已知 $J(\Q)$ 没有非平凡 $\ell$-挠点，则
\[
T_\ell(J)^{G_\Q}=0.
\]
\end{problem}
\begin{solution}
一个 Tate 模元素是相容系统
\[
(P_n)_{n\ge 1},
\qquad P_n\in J[\ell^n],
\qquad \ell P_{n+1}=P_n.
\]
若它在 $G_\Q$ 下不变，那么每个 $P_n$ 都在 Galois 作用下不变，因此都定义在 $\Q$ 上，也就是
\[
P_n\in J(\Q)[\ell^n].
\]
若 Tate 模中有非零不变向量，则至少某个 $P_n$ 非零，事实上每个层次都有一个相容的有理 $\ell^n$-挠点。

反过来，若 $J(\Q)$ 连非零 $\ell$-挠点都没有，那么显然不可能存在这样的相容系统，所以
\[
T_\ell(J)^{G_\Q}=0.
\]
这是从 Tate 模固定向量到有理挠点的最基本联系。
\end{solution}

\begin{problem}{Weil 配对与辛基}
\hard
设 $J$ 是维数 $g$ 的 Abel 簇。利用完美交替配对
\[
T_\ell(J)\times T_\ell(J)\to \Z_\ell(1)
\]
说明可以为 $T_\ell(J)$ 选一组基
\[
e_1,\dots,e_g,f_1,\dots,f_g
\]
使配对矩阵具有标准辛形状。
\end{problem}
\begin{solution}
这是带完美交替配对的自由 $\Z_\ell$-模的标准结构定理。因为配对非退化，可以先取一个原始向量 $e_1$，再选 $f_1$ 使
\[
\langle e_1,f_1\rangle=1
\]
在适当识别下成立。随后考虑它们的正交补
\[
\{x:\langle x,e_1\rangle=\langle x,f_1\rangle=0\}.
\]
交替性与完美性保证这个正交补仍是自由模，并且秩下降 $2$。

对该正交补重复同样过程，归纳下去，就得到一组基
\[
e_1,\dots,e_g,f_1,\dots,f_g
\]
满足
\[
\langle e_i,e_j\rangle=0,
\qquad
\langle f_i,f_j\rangle=0,
\qquad
\langle e_i,f_j\rangle=\delta_{ij}.
\]
因此配对矩阵正是标准辛矩阵。这也再次说明 Tate 模的秩必须是 $2g$。
\end{solution}

\begin{problem}{有理新形式给出椭圆曲线商}
\hard
设 $f$ 是权 $2$ 新形式，且 $K_f=\Q$。证明 $A_f$ 是椭圆曲线；若再知道 $S_2(\Gamma_0(37))$ 有两条这样的新形式，解释为什么 $J_0(37)$ 有两个椭圆曲线商。
\end{problem}
\begin{solution}
由
\[
\dim A_f=[K_f:\Q]
\]
立刻得到
\[
\dim A_f=1.
\]
而一维 Abel 簇就是椭圆曲线，所以 $A_f$ 是椭圆曲线。

若 $S_2(\Gamma_0(37))$ 中有两条有理新形式 $f_1,f_2$，那么对应得到两个一维模 Abel 簇 $A_{f_1},A_{f_2}$。Hecke 分解说明 $J_0(37)$ 的不同新形式方向分别给出不同的商，因此存在两个椭圆曲线商，并且
\[
J_0(37)\sim A_{f_1}\times A_{f_2}.
\]
这就是二维 Jacobian 中出现两个模椭圆曲线因子的原因。
\end{solution}

\begin{problem}{从 Jacobian 到初步 Galois 分解}
\hard
说明为什么一旦有同种分解
\[
J_0(N)\sim \prod_f A_f,
\]
就会诱导出
\[
T_\ell(J_0(N))\otimes \Q_\ell \cong \bigoplus_f \bigl(T_\ell(A_f)\otimes \Q_\ell\bigr),
\]
并解释每个 $f$-方向的维数为什么是 $2[K_f:\Q]$。
\end{problem}
\begin{solution}
同种在 Tate 模上会变成同构（至少在张量 $\Q_\ell$ 后），因为有限核与有限余核在 $\Q_\ell$-向量空间里都会消失。因此若
\[
J_0(N)\sim \prod_f A_f,
\]
便有
\[
T_\ell(J_0(N))\otimes \Q_\ell \cong \bigoplus_f \bigl(T_\ell(A_f)\otimes \Q_\ell\bigr).
\]

现在若 $\dim A_f=[K_f:\Q]$，则 Tate 模的 $\Z_\ell$-秩是它的两倍，所以
\[
\dim_{\Q_\ell}\bigl(T_\ell(A_f)\otimes \Q_\ell\bigr)=2\dim A_f=2[K_f:\Q].
\]
因此 Jacobian 的 $\ell$-进线性空间按新形式方向分块后，每一块的大小正由系数域次数控制。这就是第 9 章更深入的 Galois 表示分解的几何起点。
\end{solution}
